\chapter{Hematuria}

\section*{Definition}
Hematuria is the presence of blood in the urine, classified as gross (visible to the naked eye, appearing pink, red, or cola-colored) or microscopic (detected only on urinalysis or microscopy, defined as three or more red blood cells per high-power field). It may be symptomatic or asymptomatic and can originate from any site along the urinary tract from the glomerulus to the urethra.

\section*{Pathophysiology}
Hematuria results from disruption of vascular integrity within the urinary tract due to inflammation, infection, calculi, trauma, neoplasm, or glomerular basement membrane damage. Glomerular hematuria (dysmorphic RBCs, RBC casts) arises from immune complex deposition, basement membrane thinning, or podocyte injury causing filtration barrier disruption. Nonglomerular hematuria (eumorphic RBCs) results from mucosal bleeding due to stones, malignancy, infection, or trauma along the collecting system.

\section*{Causes}
\begin{itemize}
  \item Urinary tract infection (cystitis, pyelonephritis, prostatitis, urethritis)
  \item Nephrolithiasis (kidney stones, ureteral stones, bladder stones)
  \item Benign prostatic hyperplasia (BPH) — common in older men
  \item Glomerulonephritis (IgA nephropathy, post-streptococcal, membranous, crescentic)
  \item Malignancy (renal cell carcinoma, urothelial carcinoma of bladder/ureter, prostate cancer)
  \item Trauma (straddle injury, flank trauma, catheter-related)
  \item Strenuous exercise (benign, resolves with rest)
  \item Anticoagulant use (warfarin, DOACs, heparin) — often exacerbates underlying pathology
  \item Polycystic kidney disease
  \item Sickle cell disease or trait (papillary necrosis, hematuria)
\end{itemize}

\section*{When to See a Doctor}
See your doctor if:
\begin{itemize}
  \item Gross hematuria (visible blood in urine) regardless of associated symptoms
  \item Microscopic hematuria persisting on repeat testing
  \item Hematuria accompanied by flank pain, fever, or difficulty urinating
  \item Painless hematuria in older adults (must exclude malignancy)
  \item Hematuria in people with smoking history or occupational exposure to carcinogens
\end{itemize}

\section*{Related Symptoms}
\begin{itemize}
  \item Dysuria (painful urination) suggesting infection
  \item Flank pain radiating to groin suggesting ureteral stone
  \item Frequency, urgency, and suprapubic discomfort
  \item Passage of clots (irregular shapes suggest bladder source)
  \item Foamy urine (suggests proteinuria with glomerular disease)
  \item Lower extremity edema or hypertension (suggests glomerulonephritis)
\end{itemize}

\section*{Self-Care / Home Management}
\begin{itemize}
  \item Drink adequate water throughout the day unless fluid restriction is medically indicated.
  \item Avoid bladder irritants including caffeine, alcohol, acidic foods, and artificial sweeteners.
  \item Practice good hygiene and urinate promptly when the urge arises to prevent urinary stasis.
  \item Apply a heating pad to the lower abdomen or back for comfort during urinary discomfort.
  \item Seek medical attention for fever, flank pain, visible blood in urine, or difficulty urinating.
\end{itemize}

\section*{How Your Doctor Will Evaluate You}
\begin{itemize}
  \item History covers urinary symptoms (frequency, urgency, dysuria, hesitancy, stream changes), fever, and flank pain.
  \item Urinalysis with microscopy and culture is the first-line diagnostic test; urine cytology is obtained when malignancy is suspected.
  \item Imaging (CT urogram, renal ultrasound) evaluates for stones, obstruction, structural abnormalities, or masses.
  \item Cystoscopy is indicated for persistent hematuria or when bladder pathology is suspected, particularly in people with risk factors for bladder cancer.
\end{itemize}

\section*{Treatment Options}
\begin{itemize}
  \item Urinary tract infections are treated with targeted antibiotics based on culture sensitivity results.
  \item Nephrolithiasis management includes hydration, analgesics, medical expulsive therapy (alpha-blockers), or urologic intervention (ESWL, ureteroscopy) for large stones.
  \item Glomerular disease requires nephrology consultation for immunosuppression, blood pressure control, and proteinuria reduction.
  \item Urology referral is indicated for recurrent hematuria, suspected malignancy, obstructive symptoms, or definitive management of stones.
\end{itemize}

\clearpage
